\begin{theorem}[S-34: Determinant descent excludes every rank-one equivariant balanced seed]
\label{thm:semireg-rm-rank-one-equivariant}
Let $(X,\theta)$ be a polarized complex abelian fourfold
with Rosati-invariant real multiplication by a genuine
real quadratic field $F$, of type $(2,2)$. Write
$\theta=\theta_1+\theta_2$ and
\[
 v_q=\left(1-\frac q2\theta_1^2\right)
       \left(1-\frac q2\theta_2^2\right),\qquad q\in\mathbf Q_{>0}.
 \tag{S34.1}
\]
Let $G\subset X\times\widehat X$ be a finite effective
subgroup with a specified coherent action on
$\operatorname{Perf}(X)$ by translation and tensor product
with degree-zero line bundles. Let $E$ be a perfect complex
of virtual rank one satisfying
$e^B\operatorname{ch}(E)=v_q$ for a rational algebraic
divisor class $B$. No lift $\widetilde E$ to
$\operatorname{Perf}(X)^G$ which is simple and gluable
in that category is fully semiregular.

Here simplicity means
$\operatorname{End}_{\operatorname{Perf}(X)^G}
(\widetilde E)=\mathbf C$, and gluability means the
vanishing of its negative self-Ext groups. Neither condition
is required of the underlying object $E$. No maximal-order
or minimal-divisor-space hypothesis is imposed.
\end{theorem}
\begin{proof}
Write the action functors and their coherent multiplication
as
\[
 \phi_g(K)=t_{a_g}^*K\otimes P_g,
 \qquad
 \mu_{g,h}:\phi_g\phi_h\xrightarrow{\sim}\phi_{gh}.
\]
A lift consists of maps
$\lambda_g:\phi_g(E)\xrightarrow{\sim}E$ satisfying
\[
 \lambda_{gh}\circ\mu_{g,h,E}
       =\lambda_g\circ\phi_g(\lambda_h).
 \tag{S34.2}
\]
All natural identifications below use these specified
multiplication maps; replacing them by arbitrary
isomorphisms of the same functors would lose the argument.

Put $L=\det(E)$. The determinant of a perfect complex and
its compatibility with pullback are supplied by
Knudsen--Mumford, \emph{The projectivity of the moduli
space of stable curves I}, Theorem~2, p.~42, with the
extended determinant construction on pp.~40--42.
For a line bundle $P$ and a perfect complex of virtual
rank $r$, the local alternating determinant construction
gives
\[
 \det(E\otimes P)=\det(E)\otimes P^r,
 \qquad \det(c\operatorname{id}_E)=c^r\operatorname{id}_{\det E}.
 \tag{S34.3}
\]
Since $r=1$, taking determinants in (S34.2) gives maps
\[
 \delta_g:t_{a_g}^*L\otimes P_g\xrightarrow{\sim}L
 \tag{S34.4}
\]
with exactly the original scalar coherence, not a nontrivial
power of it. More explicitly, composition of (S34.4) for
$h$ and then $g$ agrees with (S34.4) for $gh$ after the
line-bundle identification induced by $\mu_{g,h}$.
This follows by applying the determinant functor to
(S34.2); the scalar part of $\mu_{g,h}$ is unchanged by
(S34.3). The formula can also be checked on any local
bounded complex of vector bundles, where the exponent
is the alternating sum of its ranks, namely one.

Conjugate the action by $K\mapsto K\otimes L^{-1}$.
Its functor for $g$ becomes
\[
 K\longmapsto t_{a_g}^*K\otimes R_g,
 \qquad
 R_g=t_{a_g}^*L\otimes P_g\otimes L^{-1}.
 \tag{S34.5}
\]
The maps $\delta_g$ trivialize $R_g$, and the coherence
just proved says that these trivializations identify the
entire action with the ordinary coherent translation action
of $G$ through $g\mapsto a_g$. They cancel any scalar
two-cocycle as well as the tensor factors. In particular
$E'=E\otimes L^{-1}$ acquires an ordinary translation
linearization. If $a_g=0$, (S34.4) implies $P_g\cong
\mathcal O_X$. Effectivity then gives $g=1$. Thus
$G\to X$ is injective. Its image $H$ acts freely, and
\[
 \pi:X\longrightarrow Y=X/H
 \tag{S34.6}
\]
is a finite etale isogeny of degree $|G|$.

For clarity, perfect descent here can be realized directly,
without assuming that descent in an arbitrary derived
homotopy category is effective. On $\pi_*E'$ the ordinary
linearization gives a genuine finite-group action. Form its
averaging idempotent and its image
\[
 E_Y=\operatorname{im}\left(
  \frac1{|G|}\sum_{g\in G}g:\pi_*E'\longrightarrow\pi_*E'
                         \right).
 \tag{S34.7}
\]
The pushforward is perfect: etale locally on $Y$ the
cover is a disjoint union of copies of $Y$, and the
assertion is then the preservation of perfection by a
finite direct sum. Perfection is fpqc local by Stacks
Project, Lemma~36.12.4, tag \texttt{09UG}.
The perfect category is closed under direct summands, so
the idempotent has a perfect image. After pulling back to
$X$, flat base change identifies $\pi^*\pi_*E'$ with
$\bigoplus_{h\in H}t_h^*E'$. The linearization identifies
the invariant summand with the diagonal copy of $E'$.
Consequently $\pi^*E_Y\cong E'$, compatibly with descent.
The same construction on morphism complexes, or the
adjunction and exactness of finite-group invariants, gives
for every integer $j$
\[
 \operatorname{Ext}^j_Y(E_Y,E_Y)
      \cong \operatorname{Ext}^j_X(E',E')^G
      \cong
      \operatorname{Ext}^j_{\operatorname{Perf}(X)^G}
                   (\widetilde E,\widetilde E).
 \tag{S34.8}
\]
Thus the descended object is simple and gluable. The
underlying sheaf descent used by this construction is
Stacks Project, Proposition~35.5.2, tag \texttt{023T};
the explicit finite averaging construction supplies the
perfect-complex statement and the Ext comparison.

Full semiregularity is also preserved in this descent.
One may use the enhanced equivalence
$\operatorname{Perf}(Y)\simeq
\operatorname{Perf}(X)^H$ just constructed and naturality
of the full Hochschild trace, as in S-10 and Perry 2026,
Remark~2.3. Equivalently, suppose
$u\in\operatorname{Ext}^2_Y(E_Y,E_Y)$ has zero full trace.
Its pullback has zero full trace on $X$. Under (S34.8)
it is an invariant class. The forgetting map on
$HH_{-2}$ is injective for this translation action by
the identity-sector argument of S-21. Naturality
therefore makes the equivariant full trace of $u$ zero;
full semiregularity of the lift forces $u=0$.
Tensoring by $L^{-1}$ is an equivalence and preserves
this property. Hence a fully semiregular lift would
give a fully semiregular rank-one perfect $E_Y$ on $Y$.

The character and the real-multiplication data descend
with the required normalization. The degree-two component
of $e^B\operatorname{ch}(E)=v_q$ gives $B=-c_1(E)$,
so $\operatorname{ch}(E')=v_q$.
Conjugating the rational $F$-action by the isogeny
defines an $F$-action on $Y$, and
$\theta_Y=(\pi^*)^{-1}\theta$ is a rational ample class
for which that action is Rosati invariant. Choose a
positive integer $m$ such that $\theta'=m\theta_Y$
is an integral polarization. With the corresponding
two real components of $\theta'$, the equality of
characters reads
\[
 \operatorname{ch}(E_Y)
   =\left(1-\frac{q/m^2}{2}(\theta'_1)^2\right)
      \left(1-\frac{q/m^2}{2}(\theta'_2)^2\right).
 \tag{S34.9}
\]
All equalities follow by pulling back along $\pi$,
which is an isomorphism on rational cohomology. The
polarization volume changes by
$\int_Y(\theta')^4/24=m^4D/|G|$, and $q/m^2>0$.
S-30 excludes a simple gluable fully semiregular
rank-one object with (S34.9), independently of the
polarization degree or an integral maximal-order action.
This is the desired contradiction.
\end{proof}
